% Appendix J -- What Did Not Hold (body). Sources: main/reports/FINAL_A.md 4.2-4.3, rss_ext/CLAIMS_rss.md 3,
%   tmlr_ext/CLAIMS_tmlr.md 3; outcomes from Table tab:tests (keys b.*, e11.*, e17.*, dw.*, x17.*).
\textbf{Tests and predictions that did not hold.}
\begin{itemize}\setlength{\itemsep}{1pt}\setlength{\parskip}{0pt}
\item \emph{The OEI tests.} The original mechanism tests failed: \shiter's OEI was not near zero (H2a, \nv{b.h2a.n} of 4 mechanism tasks; push \nv{b.h2a.push.shiter}), also when the executed plan was re-scored with an independent draw (\nv{b.h2aindep.n} of 4), and the OEI did not rank the damage better than the one- and three-step errors (H2b: \nv{b.h2b.td} in TD-MPC2, \nv{b.h2b.dw} in DINO-WM, against a threshold of 0.15). The first failure is the winner's curse that the theory predicted (Appendix~\ref{app:theory}); after E1 the mechanism claim was narrowed to drift, tested on new seeds for TD-MPC2 and on the E1 cell for PushT (H2a$'$).
\item \emph{Near-tie density.} The benefit of sharing a draw did not grow with the density of near ties (H4b: Spearman \nv{b.h4b.rho}), although sharing helped in most tasks.
\item \emph{DMC.} The variant \shiter repaired \nv{e1.dmc.nfix} of \nv{e1.dmc.naff} damaged DMC cells in E1 (H3a fails for DMC) and was harmful on walker at 3 bits ($D_{\shiter}-D_{\dtm}$ \nv{b.h3a.walker}\% of return); in E11 it repaired \nv{e11.dmc.nfix.shiter} of \nv{e11.dmc.ndam}. With 12 iterations the better dither repairs \nv{e15.it1.n} of the 4 tasks (It1), but walker stays unrepaired even at 24 iterations ($\red$ \nv{e15.walker.red24}).
\item \emph{Better post-training quantizers.} On top of the best post-training quantizer, the extra gain of \shiter was borderline in TD-MPC2 (\nv{b.h3b.td} cells, exactly the threshold) and found in DINO-WM in only one of three cells (\nv{b.h3b.dw}), where the quantizer alone left little damage.
\item \emph{The PushT positive control.} Holding the full-size shared dither fixed for the plan (\shplan) should have brought back \dtm's drift on PushT, as it did in TD-MPC2; its drift ratio was \nv{e2.drift.pusht.shplan}. Only coarse tables were exploited in DINO-WM (Appendix~\ref{app:mech}). H2c also failed on PushT (Spearman of $N$ and damage \nv{b.h2c.pusht.d}), and the E3 expectation that tables kept per candidate behave like redrawn per-candidate dither failed in \nv{e.e3.indstep.n}--\nv{e.e3.indep.n} of 12 TD-MPC2 tasks.
\item \emph{Wall in E1.} None of the planned Wall cells was damaged, and at 4 bits \dtm beat the unquantized model, so Wall contributed only an exploratory cell to E1.
\item \emph{Fresh noise (B3).} We predicted that adding fresh noise to \dtm would not repair a majority of the damaged cells in any benchmark; it repaired \nv{e11.mw.nfix.detnoise} of \nv{e11.mw.ndam} Meta-World cells, though only \nv{e11.dmc.nfix.detnoise} of \nv{e11.dmc.ndam} DMC cells and neither damaged DINO-WM cell, and on the bit axis it held up less well than dither (Section~\ref{sec:noise}).
\item \emph{Dg1 on DINO-WM.} The summary statistic did not beat the latent one-step error by the planned margin on DINO-WM's single damaged cell (\nv{dg1.dw}).
\item \emph{PushT at 4 bits on the new split.} The E1 damage (\nv{e1.pusht4.D.det} points) did not reappear on the new split (\nv{dw.pusht4.D.det}); re-running the E1 pairs on the new GPU type reproduced it (\nv{x17.D.det}), so the split, not the hardware, explains the difference (post hoc, X17).
\item \emph{Latency and formats.} The kernel overhead of \detshift exceeded 5\% at $N=512$ (K1 held in only \nv{e17.k1.n} of \nv{e17.k1.of} conditions; \nv{e17.detshift.5m512}\% and \nv{e17.detshift.317m512}\%). Pre-converted BF16 was not lossless (K3: \nv{e17.k3.n} of \nv{e17.k3.of} tasks without damage; door-open \nv{e17.bf16.dooropen} points, cheetah-run \nv{e17.bf16.cheetah}\%). The integer rollout was not faster than pre-converted BF16 (L2c: \nv{b.l2c.range}$\times$ its time), nor was the whole planning step (S2: \nv{b.s2.range}$\times$). In eager execution the integer kernel was slower than FP32 (E7), and H5's cases where dither with fewer bits Pareto-dominates \dtm with more were weak (push only). NVFP4 caused too little damage for F1 (\nv{b.f1.nv} of 12 tasks), and on DMC the FP4 cells were not repaired (F2: \nv{b.f2.dmc}).
\item \emph{Few iterations and weights.} With one MPPI iteration \shiter is no better than \dtm, and in the E7 kernel at three iterations dither was harmful on DMC. Stochastic rounding of weights helped in TD-MPC2 and hurt in DINO-WM.
\end{itemize}

\textbf{Claims we do not make.}
\begin{itemize}\setlength{\itemsep}{1pt}\setlength{\parskip}{0pt}
\item That the OEI predicts damage better than one-step error, or that planners do not exploit redrawn noise: the winner's curse of the last selection arises under fresh noise too.
\item That one-step or multi-step error is the wrong measure in general. Our statement concerns the variants at one bit width; across bit widths the one-step error ranks the damage, errors over several steps separate \dtm in DINO-WM, and the reward error partly separates it in TD-MPC2.
\item That persistence decides the outcome in both families, that DINO-WM is dominated by mean bias (\detbc is worse than \dtm on PushT), or that the model family causes which component dominates: E16 swapped only the planners' update rules, and the families differ in more than the planner.
\item That any noise or that randomization alone removes the damage, or that a low persistent share implies repair: \detnoise repairs Meta-World but not DMC or DINO-WM.
\item That a loss of resolution causes \dtm's damage in TD-MPC2: fixed random thresholds repair at most \nv{e11.mw.nfix.shep} of \nv{e11.mw.ndam} Meta-World cells.
\item That bias correction is a remedy, that \shiter alone gains a bit in most Meta-World tasks (it does so in \nv{e12.pt.nG.shiter} per tensor and \nv{e12.pc.nG.shiter} per channel, of \nv{e12.mw.n}), or that dither gains two bits on Wall (\nv{dw.wall.G.range}).
\item That $\rho_p\cdot\sigma_e/\sigma_J$, measured before deployment, predicts the damage: it nearly restates the manipulated variable and fails on DINO-WM.
\item That \detshift comes without latency cost, that pre-converted BF16 is lossless, that our integer W8A4 world model is faster than BF16, that 4-bit activations run faster than 8-bit ones on the GPUs we used, that low-bit arithmetic speeds up the whole planning step as much as the rollouts, or that the world model dominates the cost of a planning step.
\item That the fix works on DMC at the default number of iterations, with a single planner iteration, or in every task and planner setting.
\item That dither helps in general on top of strong post-training quantization, or any conclusion about quantization-aware training, whose baseline was weak.
\item That accelerator FP4 formats reproduce the failure of uniform grids, or that larger models are more robust to rounding: the primary bit widths at 317M are close to those at 5M.
\item That every persistent error is exploited, that a fixed dither obeys a $1/N_{\mathrm{eff}}$ law, or that sharing and redrawing are both always needed: in DINO-WM the full-size fixed table was not exploited.
\item That the PushT damage at 4 bits replicates on the new split.
\end{itemize}
